\section{五矩双流体的两条可实施路线：DG与全隐式有限体积}
\subsection{DG：一个单元内保留多个自由度}
\citet{L156}将两物种Euler方程与完整Maxwell系统写成
$\partial_tU+\nabla\cdot F(U)=S(U)$，每单元用低阶多项式展开$U_h=\sum_jU_j(t)\varphi_j$。
乘试验函数、单元积分并分部积分后，有
\begin{equation}
 M_{ij}\dot U_j=\int_K\nabla\varphi_i\cdot F(U_h)\dd V
 -\int_{\partial K}\varphi_i\widehat F\dd S+\int_K\varphi_iS(U_h)\dd V.
\end{equation}
正交基使质量矩阵$M$对角；单元间通过数值通量$\widehat F$交换信息。
该文一次基用于二阶空间精度，二次基用于三阶精度；二维基分别包含$\{1,x,y\}$及增加的$xy,x^2-1/3,y^2-1/3$。
边积分分别用2或3个Gaussian点、体积分4或9个点；局部Lax通量按物种声速和流速选耗散，电磁部分含光速。

限制器也是算例定义的一部分。该文使用特征空间minmod/TVB；一旦线性系数被限幅，高阶系数清零，因此“采用三阶DG”不意味着每个单元始终三阶。
时间采用SSPRK3统一推进通量与刚性源项。对纯振荡$\dot q=i\omega q$，二阶显式RK的放大因子模大于1，不能以流体平流CFL代替源项稳定性分析。
文中典型要求$\Delta t<\omega_{pe}^{-1}$，还受光波CFL约束。所读预印本第12页三阶CFL的$c/\Delta x$量纲错误；应具有$\Delta x/c$量纲，具体安全系数随基阶和空间维数而变。

GEM部分使用$m_i/m_e=25$、$c/\VA=10$、$T_e/T_i=0.2$，比较$128\times64$与$512\times256$网格，时间步分别$0.002$和$0.0005\Omega_{ci}^{-1}$。
粗格二阶持续出岛而三阶没有同样的岛；细格两阶通量接近。更强限制/耗散也可能消岛，说明不能给数值耗散与磁岛数建立简单单调关系。
晚期局地动量对浮点精度有约10\%差别，且磁岛偏移方向可受舍入误差影响。
这要求在物理结论旁报告网格、阶数、限制器、时间步和误差敏感性，而不只列网格数。

\subsection{隐式有限体积：稳定步长、矩阵内存与求解成本}
\citet{L243}采用18变量五矩双流体--Maxwell--散度清理系统。
二阶空间重构使用最小二乘梯度及Venkatakrishnan限制器，物种通量用AUSM+up，电磁通量用改进Rusanov；隐式BDF2配Newton迭代，数值差分组装Jacobian，再以GMRES/Jacobi求解。
它实际组装矩阵，不能笼统归入无矩阵Newton--Krylov算法。

GPU实现以一个线程处理一个单元，计算通量、源项与变量扰动后的Jacobian块；CSR矩阵装配移到设备侧，避免反复CPU传输。
边界处理等仍有主机工作。波算例检查二阶误差；GEM采用双电流片全周期边界、$256^2$与$512^2$单元、质量比25，至$t\Omega_{ci}=25$共10000步。
后者$\Delta t\Omega_{ci}=0.0025$，P100完整一次计算约49小时16分；没有因“全隐式”就无限加大物理时间步。

报告的14.23倍加速相对同机\emph{单个CPU核心}；若换更快的K40主机CPU作基准，约为8.15倍。
小问题GPU占用低，某些PARALUTION配置比CPU慢；大问题又受到完整耦合矩阵显存限制。这些是当时PETSc 3.9、PARALUTION 1.1与特定硬件的结果，不能用于预测今日任意平台。
GPU/CPU磁通吻合验证实现一致性，尚不验证电子压力闭合或局地非理想电场的物理准确性。

\section{线性、非线性与湍流算例的实际分辨要求}
\begin{longtable}{P{29mm}P{64mm}P{55mm}}
\toprule 文献与路线 & 已核实的计算细节 & 必须保留的限制\\\midrule\endhead
\citet{L224}：斜向惯性撕裂 & HiFi谱元/BDF2；旋转二维平面模拟不同倾角；$m_i/m_e=100$，$d_e$约8--11格，导场比1--10，超电阻扫描 & 各角度独立二维线性计算，不等于一次三维非线性竞争\\
\citet{L268}：电阻Hall本征值 & 不可压Harris平衡；SciPy四阶配点，横向区间约$[-15,15]$，同时扫描两方向波数 & 是本征边值问题，未列统一固定网格/容差；增长与传播频率应分别核对\\
\citet{L191,L214,F697}：爆发阶段 & 早期$x$约$10^4$差分点、$y$约200模；后续二维谱/RK4达到$4096^2$--$8192^2$ & 增加横向模数后看到次级岛，旧的平滑加速结论受分辨率影响\\
\citet{E471,E472}：回旋流体 & 通常$1024\times128$，部分非线性提高到256横向模；Padé回旋算子 & 压力/导心密度能不同于热化；总能误差约百分数不能替代热预算\\
\citet{E474}：有限$\beta_e$ & 谱空间、三阶Adams--Bashforth、短尺度滤波，最高$2880^2$，能量分析限于损失小于1\%时段 & 滤波后段不宜推断物理耗散；移动X点影响增长率曲线\\
\citet{F644}：KH卷涡 & 2.5维Hall，$4096\times8192$，盒约$400\times402d_i$，RK3，滤除$kd_i\gtrsim3$ & 无显式电阻、惯性及张量破冻结项；拓扑改变受滤波/离散影响\\
\citet{L294}：非均匀片传播 & F3D，$102.4\times51.2\times256d_i$，网格间距$(0.05,0.05,1)d_i$，质量比25 & 前沿判据$B_y=0.04$与非均匀片几何是拟合的一部分\\
\bottomrule
\end{longtable}

\subsection{嵌入Hall区时，过渡函数本身会进入方程}
\citet{L035}将局部Hall计算嵌入外部MHD，比较显示磁通演化接近全域Hall。
但若直接把非理想电场写成$\alpha(\vct x)\E_H$，感应方程中出现
\begin{equation}
 -\nabla\times(\alpha\E_H)=-\alpha\nabla\times\E_H-\nabla\alpha\times\E_H.
\end{equation}
第二项是过渡带带来的额外作用。论文给出的平滑erf混合不等于已证明守恒的多尺度耦合算法。
未来复现应同时改变Hall区域大小和过渡带宽度，检查过渡区能量及电流，而不只比较一个全局通量曲线。

\subsection{三种“率”应在数据文件里就分开}
保存$E_z/(B_mV_{\!A,m})$、$\dd\ln\psi/\dd t$和$\delta/L$三个独立量。
\citet{L264}的重联频率由磁通演化时标获得；\citet{F644}的0.1--0.3主要是电流片纵横比估计。
只有满足稳态、不可压、明确的出流速度和上游归一化时，纵横比才与电场率相联系。
在合并、湍流和非稳态小片中，这些量即使数值接近也不可互换。
